\documentclass{situationist-preprint}

\title{Aboutness Is Primitive}
\msid{MND-12-a}
\msfamily{MND-12}
\msdiscipline{Philosophy of Mind}
\mskeywords{intentionality; naturalization; teleosemantics; indeterminacy; content determinacy; Swampman}
\mstype{Principal manuscript}

\begin{document}
\maketitle

\begin{abstract}
Can intentionality be naturalized, that is, can the aboutness of mental states be analyzed in terms of causation, information and biological function? We argue that it cannot: aboutness is primitive. Every naturalistic condition on content is fixed by a finite base of facts about a creature's history, dispositions and environment, and every such base is compatible with distinct content assignments that the creature's states nonetheless discriminate. The indeterminacy problems that have defeated causal, informational and teleological theories one by one are instances of this structural fact. But primitiveness of aboutness does not leave its distribution unexplained. Teleosemantics explains which states have content and the range within which their contents fall: a state is about something within the range of conditions it was selected or trained to track. Aboutness is thus a relation that natural theory constrains but does not define. The view is compatible with physicalism, since primitiveness concerns analysis and explanation, not fundamentality; it accounts for content in artificial systems; and it gives a principled verdict on Swampman.
\end{abstract}

\section{The problem}

\citet{Brentano1874} took intentionality, the directedness of mental states on objects, to be the mark of the mental, and held that no physical phenomenon exhibits it. The naturalistic program of the late twentieth century set out to show that he was wrong, by giving necessary and sufficient conditions for a state's being about something in terms that do not themselves presuppose aboutness. \citet{Fodor1987} stated the motivation: if aboutness is real, it must be something else.

Three families of theory were proposed. Informational theories take a state to be about what it carries information about \citep{Dretske1981}; they face the problem of misrepresentation, since a state that is caused by cows and by horses on dark nights carries information about the disjunction. Fodor's asymmetric dependence theory takes a symbol to mean cow if horse-caused tokenings depend on cow-caused ones but not conversely \citep{Fodor1990}; it faces counterexamples in which the dependences run the wrong way or cannot be evaluated. Teleosemantic theories take a state to be about what it has the biological function of indicating, or what must obtain for its consumers to perform their functions \citep{Millikan1984,Millikan1989,Papineau1987}; they solve the misrepresentation problem, since a state can fail to perform its function, but face the indeterminacy problem: the frog's fly-detector was selected for responding to flies, to frog food, and to small dark moving things, and the selection history does not choose among these \citep{Neander2017}. \citet{Davidson1987} added the Swampman case: a molecule-for-molecule duplicate of a person, formed by chance, has no history, and so on teleosemantics no contentful states.

Underlying these specific failures is a more general argument. \citet{Kripke1982}, developing Wittgenstein, argued that no fact about a speaker's past behavior and dispositions determines whether she means addition or a function that agrees with addition on the cases she has computed; \citet{Quine1960} had argued for an indeterminacy of translation from the behavioral evidence. Naturalists have replied by appeal to the naturalness of the candidate contents \citep{Lewis1983}, and by refining teleosemantics with causal and similarity constraints \citep{Neander2017,Shea2018}. Opponents of naturalization have proposed that intentionality is grounded in phenomenal consciousness \citep{Mendelovici2018}.

We argue for a position between these: aboutness cannot be naturalized, and is primitive; but its distribution, which states have content and roughly what they are about, is explained by teleosemantics. The view is the one MND-07-a appeals to when it holds that symbols in a suitably embedded system can have content.

\section{Framework}

\begin{definition}[Naturalization]\label{d:nat}
Intentionality is \emph{naturalized} if there is a condition $N$, stated in non-intentional terms (causal, informational, nomological or teleological), such that for every possible state $s$ and content $c$, $s$ is about $c$ if and only if $N(s,c)$, and $s$ is about $c$ because $N(s,c)$.
\end{definition}

\begin{definition}[Primitive]\label{d:prim}
Aboutness is \emph{primitive} if intentionality is not naturalized: no non-intentional condition both fixes and explains which contents states have.
\end{definition}

\begin{definition}[Distribution]\label{d:dist}
The \emph{distribution} of aboutness is the pattern of facts about which states have contents and within what range their contents fall. A theory \emph{explains the distribution} if it says why, among the states of creatures of a given kind, those states and not others are contentful, and why their contents lie within the ranges they do.
\end{definition}

\begin{principle}[Finite base]\label{p:base}
For any naturalistic condition $N$ and any finite creature, whether $N(s,c)$ holds is fixed by a finite base of facts: the creature's history of selection and learning, its actual and counterfactual dispositions within the range of circumstances it could encounter, and the causal relations between its states and its environment.
\end{principle}

\begin{principle}[Underdetermination]\label{p:under}
For any finite base, there are distinct contents $c$ and $c'$ that agree on every fact in the base: they are satisfied in the same historical, dispositional and causal circumstances.
\end{principle}

\begin{principle}[Determinacy]\label{p:det}
For some states and some pairs of contents $c$ and $c'$ that agree on every fact in the relevant base, the state is about $c$ and not about $c'$.
\end{principle}

Principle~\ref{p:base} expresses what makes a condition naturalistic: it is fixed by the kinds of fact natural science records about an organism and its surroundings. Principle~\ref{p:under} generalizes the indeterminacy cases. Contents are individuated more finely than any finite pattern of causes, dispositions and selective successes: fly and small dark moving thing agree on the frog's history; addition and quaddition agree on every computation a finite subject could perform; rabbit and undetached rabbit part agree on every stimulus. Principle~\ref{p:det} is the datum that our thoughts are determinate in some such cases. When I think that two and two make four I think of addition, not of quaddition; when I think of a rabbit I do not think of an undetached rabbit part. The datum is introspective in part, but it is also required by the practice of reasoning: inference preserves truth only if premises and conclusions have the contents they have.

\section{The argument}

\begin{lemma}\label{l:fail}
No naturalistic condition fixes content.
\end{lemma}

\begin{proof}
Let $N$ be a naturalistic condition. By Principle~\ref{p:det}, there is a state $s$ and contents $c$, $c'$ agreeing on the base relevant to $s$ such that $s$ is about $c$ and not about $c'$. By Principle~\ref{p:base}, whether $N(s,c)$ and whether $N(s,c')$ are fixed by that base, and since $c$ and $c'$ agree on it, $N(s,c)$ holds if and only if $N(s,c')$ does. If $N(s,c)$ holds, $N$ assigns $c'$ as well, which $s$ is not about; if it does not, $N$ fails to assign $c$, which $s$ is about. Either way $N$ does not satisfy the biconditional of Definition~\ref{d:nat}.
\end{proof}

\begin{theorem}\label{t:prim}
Intentionality cannot be naturalized; aboutness is primitive.
\end{theorem}

\begin{proof}
Naturalization requires a condition satisfying the biconditional of Definition~\ref{d:nat}. By Lemma~\ref{l:fail} there is none, and by Definition~\ref{d:prim} aboutness is primitive.
\end{proof}

The argument does not depend on the details of any theory, and it explains why the theories failed as they did. Each refinement of teleosemantics narrows the range of contents compatible with the base; none can narrow it to one in every case, because the base is finite and contents are individuated more finely.

\begin{proposition}\label{p:dist}
Teleosemantics explains the distribution of aboutness: the contentful states of creatures with histories of selection or learning are the states that have the function of tracking conditions, and their contents fall within the range of conditions whose obtaining explains the selection or reinforcement of those states.
\end{proposition}

\begin{proof}
We argue for the explanation in three steps. First, it is not an accident that the contentful states of organisms are the states that guide behavior with respect to the environment: states with no such role, such as the temperature of the liver, are not about anything. The teleosemantic theory says why: contentful states are those that were selected, or reinforced by learning \citep{Dretske1988}, because their tokening covaried with conditions that mattered to the success of the behavior they guided. Second, the range within which a state's content falls is the range of conditions whose covariation explains that selection; the frog's state is about something in the range of fly, frog food and small dark moving thing, and not about, for instance, the color of the pond. Third, the explanation is general: it predicts which states of a newly studied creature will be contentful, and roughly what they will be about, from facts about function alone, and the predictions are borne out across perception, motor control and memory \citep{Shea2018}.
\end{proof}

Theorem~\ref{t:prim} and Proposition~\ref{p:dist} fit together. Teleosemantics fixes the distribution up to the range left open by Principle~\ref{p:under}; within that range, aboutness selects. The relation is primitive in that no non-intentional condition fixes its selection, and constrained in that it falls only on states with functions and within the ranges they determine.

\begin{remark}
Primitiveness, as Definition~\ref{d:prim} states it, concerns analysis and explanation. It does not imply that aboutness is a fundamental property. It is consistent with physicalism as MND-02 defines it, on which no fundamental property is mental: the totality of physical facts, including the intrinsic natures that MND-01 recognizes, may necessitate the facts about content without there being any naturalistic condition, in the sense of Principle~\ref{p:base}, that fixes and explains them.
\end{remark}

\section{Objections and replies}

\begin{objection}[Reference magnetism]
The more natural of two candidate contents is the one a state is about \citep{Lewis1983}. Addition is more natural than quaddition, rabbit than undetached rabbit part. Naturalness breaks the ties.
\end{objection}

\begin{reply}
Naturalness is not a condition of the kind Definition~\ref{d:nat} requires. It is a feature of the candidate contents, not of the creature's history, dispositions or causal relations, and the principle that thought gravitates toward natural contents is not stated in non-intentional terms but is a constraint on the intentional relation itself. To adopt it is to accept that aboutness has a determinant beyond the natural base, which is our thesis. Moreover, naturalness does not settle the frog's case, in which fly and small dark moving thing are both reasonably natural.
\end{reply}

\begin{objection}[Refined teleosemantics]
\citet{Neander2017} solves the indeterminacy problem for the frog: the function of the detector is to respond to small dark moving things, since that is the property to which it is causally sensitive. Further refinement can deal with the other cases.
\end{objection}

\begin{reply}
Neander's proposal selects one candidate in the frog's case, and we take no position on whether it selects correctly. It does not touch Principle~\ref{p:under}, since causal sensitivity is part of the finite base; for the content of thoughts about addition, or of concepts whose instances the subject cannot discriminate from those of an extensionally divergent concept, it leaves the alternatives tied. Refinement narrows the range, which is what Proposition~\ref{p:dist} says teleosemantics does.
\end{reply}

\begin{objection}[Swampman]
If teleosemantics explains the distribution of content, then Swampman, who has no history, has no contentful states. That is incredible.
\end{objection}

\begin{reply}
On teleosemantics as an analysis, Swampman's states lack content. On our view the teleosemantic explanation is etiological: it explains why the states of normal creatures have the contents they have, without constituting their having them. Swampman is a case in which the states have arisen by accident, and the explanation of their content is missing, but the content need not be. If the facts about content are necessitated by the totality of physical facts, as the remark above allows, Swampman's states have content, and the distribution of content in him is an accident of the same kind as the distribution of his molecules.
\end{reply}

\begin{objection}[Machines]
A primitivist view makes it mysterious how an engineered system could have contentful states, and MND-08 holds that machines can think.
\end{objection}

\begin{reply}
Proposition~\ref{p:dist} includes functions acquired by learning, and an engineered system whose states are shaped by training to track conditions has states with functions in this sense. The distribution of content in such systems is explained in the same way as in organisms. That is what MND-07-a needs: if the system in the Chinese Room is embedded so that its symbols have the function of tracking what they are about, the syntax-semantics objection does not distinguish it from a brain.
\end{reply}

\section{Conclusion}

Every naturalistic condition on content is fixed by a finite base of facts, and every finite base leaves open contents that our thoughts distinguish. No such condition can fix and explain content, and aboutness is therefore primitive. But it is not unconstrained: teleosemantics explains which states have content and within what range their contents lie. Intentionality cannot be naturalized, and its place in nature can still be described.

\begin{thebibliography}{99}\small
\bibitem[Brentano(1874)]{Brentano1874} Brentano, F. (1874). \emph{Psychologie vom empirischen Standpunkte}. Duncker \& Humblot.
\bibitem[Davidson(1987)]{Davidson1987} Davidson, D. (1987). Knowing one's own mind. \emph{Proceedings and Addresses of the American Philosophical Association}, 60(3), 441--458.
\bibitem[Dretske(1981)]{Dretske1981} Dretske, F. (1981). \emph{Knowledge and the Flow of Information}. MIT Press.
\bibitem[Dretske(1988)]{Dretske1988} Dretske, F. (1988). \emph{Explaining Behavior: Reasons in a World of Causes}. MIT Press.
\bibitem[Fodor(1987)]{Fodor1987} Fodor, J.~A. (1987). \emph{Psychosemantics}. MIT Press.
\bibitem[Fodor(1990)]{Fodor1990} Fodor, J.~A. (1990). \emph{A Theory of Content and Other Essays}. MIT Press.
\bibitem[Kripke(1982)]{Kripke1982} Kripke, S.~A. (1982). \emph{Wittgenstein on Rules and Private Language}. Harvard University Press.
\bibitem[Lewis(1983)]{Lewis1983} Lewis, D. (1983). New work for a theory of universals. \emph{Australasian Journal of Philosophy}, 61(4), 343--377.
\bibitem[Mendelovici(2018)]{Mendelovici2018} Mendelovici, A. (2018). \emph{The Phenomenal Basis of Intentionality}. Oxford University Press.
\bibitem[Millikan(1984)]{Millikan1984} Millikan, R.~G. (1984). \emph{Language, Thought, and Other Biological Categories}. MIT Press.
\bibitem[Millikan(1989)]{Millikan1989} Millikan, R.~G. (1989). Biosemantics. \emph{The Journal of Philosophy}, 86(6), 281--297.
\bibitem[Neander(2017)]{Neander2017} Neander, K. (2017). \emph{A Mark of the Mental: In Defense of Informational Teleosemantics}. MIT Press.
\bibitem[Papineau(1987)]{Papineau1987} Papineau, D. (1987). \emph{Reality and Representation}. Blackwell.
\bibitem[Quine(1960)]{Quine1960} Quine, W.~V.~O. (1960). \emph{Word and Object}. MIT Press.
\bibitem[Shea(2018)]{Shea2018} Shea, N. (2018). \emph{Representation in Cognitive Science}. Oxford University Press.
\end{thebibliography}

\end{document}
